% Vier Arten von Hüllkörpern um denselben Körper.
% Eigene Darstellung in Anlehnung an LaValle2006 (Fig. 5.9, S. 98; Beleg LaValle2006.05).
% OBB per Hauptachsen berechnet, Kugel um die Mitte des achsparallelen Quaders.
% Stil: ba-grafik.tex. Einbinden: % Vier Arten von Hüllkörpern um denselben Körper.
% Eigene Darstellung in Anlehnung an LaValle2006 (Fig. 5.9, S. 98; Beleg LaValle2006.05).
% OBB per Hauptachsen berechnet, Kugel um die Mitte des achsparallelen Quaders.
% Stil: ba-grafik.tex. Einbinden: % Vier Arten von Hüllkörpern um denselben Körper.
% Eigene Darstellung in Anlehnung an LaValle2006 (Fig. 5.9, S. 98; Beleg LaValle2006.05).
% OBB per Hauptachsen berechnet, Kugel um die Mitte des achsparallelen Quaders.
% Stil: ba-grafik.tex. Einbinden: % Vier Arten von Hüllkörpern um denselben Körper.
% Eigene Darstellung in Anlehnung an LaValle2006 (Fig. 5.9, S. 98; Beleg LaValle2006.05).
% OBB per Hauptachsen berechnet, Kugel um die Mitte des achsparallelen Quaders.
% Stil: ba-grafik.tex. Einbinden: \input{Bilder/tikz/huellkoerper} in einer figure-Umgebung.
\begin{tikzpicture}[ba]
  \newcommand{\baKoerper}{%
    \draw[ba dick, baGruen, fill=baFlaeche] (0,0) -- (1.8,0.9) -- (2.1,1.4) -- (1.7,1.6) -- (0.9,0.85) -- (0.2,0.7) -- (-0.1,0.3) -- cycle;}
  \tikzset{huelle/.style={ba normal, baBlau}}
  % a) Kugel
  \begin{scope}
    \draw[huelle] (1.0,0.8) circle (1.281);
    \baKoerper
    \node[ba teiltitel] at (1.0,-0.75) {a) Kugel};
  \end{scope}
  \draw[ba trenner] (2.75,-1.3) -- (2.75,2.2);
  % b) AABB
  \begin{scope}[xshift=3.7cm]
    \draw[huelle] (-0.1,0) rectangle (2.1,1.6);
    \baKoerper
    \node[ba teiltitel] at (1.0,-0.75) {b) AABB};
  \end{scope}
  \draw[ba trenner] (6.55,-1.3) -- (6.55,2.2);
  % c) OBB
  \begin{scope}[xshift=7.5cm]
    \draw[huelle] (1.917,1.723) -- (-0.273,0.481) -- (0.052,-0.092) -- (2.242,1.15) -- cycle;
    \baKoerper
    \node[ba teiltitel] at (1.0,-0.75) {c) OBB};
  \end{scope}
  \draw[ba trenner] (10.35,-1.3) -- (10.35,2.2);
  % d) Konvexe Hülle
  \begin{scope}[xshift=11.3cm]
    \draw[huelle] (0,0) -- (1.8,0.9) -- (2.1,1.4) -- (1.7,1.6) -- (0.2,0.7) -- (-0.1,0.3) -- cycle;
    \baKoerper
    \node[ba teiltitel] at (1.0,-0.75) {d) Konvexe Hülle};
  \end{scope}
  % Tendenz
  \draw[ba pfeil] (-0.2,-1.65) -- (13.5,-1.65);
  \node[anchor=north, align=center] at (6.65,-1.75) {von a) nach d) meist engere Anpassung, aber aufwendigerer Überlappungstest};
\end{tikzpicture}
 in einer figure-Umgebung.
\begin{tikzpicture}[ba]
  \newcommand{\baKoerper}{%
    \draw[ba dick, baGruen, fill=baFlaeche] (0,0) -- (1.8,0.9) -- (2.1,1.4) -- (1.7,1.6) -- (0.9,0.85) -- (0.2,0.7) -- (-0.1,0.3) -- cycle;}
  \tikzset{huelle/.style={ba normal, baBlau}}
  % a) Kugel
  \begin{scope}
    \draw[huelle] (1.0,0.8) circle (1.281);
    \baKoerper
    \node[ba teiltitel] at (1.0,-0.75) {a) Kugel};
  \end{scope}
  \draw[ba trenner] (2.75,-1.3) -- (2.75,2.2);
  % b) AABB
  \begin{scope}[xshift=3.7cm]
    \draw[huelle] (-0.1,0) rectangle (2.1,1.6);
    \baKoerper
    \node[ba teiltitel] at (1.0,-0.75) {b) AABB};
  \end{scope}
  \draw[ba trenner] (6.55,-1.3) -- (6.55,2.2);
  % c) OBB
  \begin{scope}[xshift=7.5cm]
    \draw[huelle] (1.917,1.723) -- (-0.273,0.481) -- (0.052,-0.092) -- (2.242,1.15) -- cycle;
    \baKoerper
    \node[ba teiltitel] at (1.0,-0.75) {c) OBB};
  \end{scope}
  \draw[ba trenner] (10.35,-1.3) -- (10.35,2.2);
  % d) Konvexe Hülle
  \begin{scope}[xshift=11.3cm]
    \draw[huelle] (0,0) -- (1.8,0.9) -- (2.1,1.4) -- (1.7,1.6) -- (0.2,0.7) -- (-0.1,0.3) -- cycle;
    \baKoerper
    \node[ba teiltitel] at (1.0,-0.75) {d) Konvexe Hülle};
  \end{scope}
  % Tendenz
  \draw[ba pfeil] (-0.2,-1.65) -- (13.5,-1.65);
  \node[anchor=north, align=center] at (6.65,-1.75) {von a) nach d) meist engere Anpassung, aber aufwendigerer Überlappungstest};
\end{tikzpicture}
 in einer figure-Umgebung.
\begin{tikzpicture}[ba]
  \newcommand{\baKoerper}{%
    \draw[ba dick, baGruen, fill=baFlaeche] (0,0) -- (1.8,0.9) -- (2.1,1.4) -- (1.7,1.6) -- (0.9,0.85) -- (0.2,0.7) -- (-0.1,0.3) -- cycle;}
  \tikzset{huelle/.style={ba normal, baBlau}}
  % a) Kugel
  \begin{scope}
    \draw[huelle] (1.0,0.8) circle (1.281);
    \baKoerper
    \node[ba teiltitel] at (1.0,-0.75) {a) Kugel};
  \end{scope}
  \draw[ba trenner] (2.75,-1.3) -- (2.75,2.2);
  % b) AABB
  \begin{scope}[xshift=3.7cm]
    \draw[huelle] (-0.1,0) rectangle (2.1,1.6);
    \baKoerper
    \node[ba teiltitel] at (1.0,-0.75) {b) AABB};
  \end{scope}
  \draw[ba trenner] (6.55,-1.3) -- (6.55,2.2);
  % c) OBB
  \begin{scope}[xshift=7.5cm]
    \draw[huelle] (1.917,1.723) -- (-0.273,0.481) -- (0.052,-0.092) -- (2.242,1.15) -- cycle;
    \baKoerper
    \node[ba teiltitel] at (1.0,-0.75) {c) OBB};
  \end{scope}
  \draw[ba trenner] (10.35,-1.3) -- (10.35,2.2);
  % d) Konvexe Hülle
  \begin{scope}[xshift=11.3cm]
    \draw[huelle] (0,0) -- (1.8,0.9) -- (2.1,1.4) -- (1.7,1.6) -- (0.2,0.7) -- (-0.1,0.3) -- cycle;
    \baKoerper
    \node[ba teiltitel] at (1.0,-0.75) {d) Konvexe Hülle};
  \end{scope}
  % Tendenz
  \draw[ba pfeil] (-0.2,-1.65) -- (13.5,-1.65);
  \node[anchor=north, align=center] at (6.65,-1.75) {von a) nach d) meist engere Anpassung, aber aufwendigerer Überlappungstest};
\end{tikzpicture}
 in einer figure-Umgebung.
\begin{tikzpicture}[ba]
  \newcommand{\baKoerper}{%
    \draw[ba dick, baGruen, fill=baFlaeche] (0,0) -- (1.8,0.9) -- (2.1,1.4) -- (1.7,1.6) -- (0.9,0.85) -- (0.2,0.7) -- (-0.1,0.3) -- cycle;}
  \tikzset{huelle/.style={ba normal, baBlau}}
  % a) Kugel
  \begin{scope}
    \draw[huelle] (1.0,0.8) circle (1.281);
    \baKoerper
    \node[ba teiltitel] at (1.0,-0.75) {a) Kugel};
  \end{scope}
  \draw[ba trenner] (2.75,-1.3) -- (2.75,2.2);
  % b) AABB
  \begin{scope}[xshift=3.7cm]
    \draw[huelle] (-0.1,0) rectangle (2.1,1.6);
    \baKoerper
    \node[ba teiltitel] at (1.0,-0.75) {b) AABB};
  \end{scope}
  \draw[ba trenner] (6.55,-1.3) -- (6.55,2.2);
  % c) OBB
  \begin{scope}[xshift=7.5cm]
    \draw[huelle] (1.917,1.723) -- (-0.273,0.481) -- (0.052,-0.092) -- (2.242,1.15) -- cycle;
    \baKoerper
    \node[ba teiltitel] at (1.0,-0.75) {c) OBB};
  \end{scope}
  \draw[ba trenner] (10.35,-1.3) -- (10.35,2.2);
  % d) Konvexe Hülle
  \begin{scope}[xshift=11.3cm]
    \draw[huelle] (0,0) -- (1.8,0.9) -- (2.1,1.4) -- (1.7,1.6) -- (0.2,0.7) -- (-0.1,0.3) -- cycle;
    \baKoerper
    \node[ba teiltitel] at (1.0,-0.75) {d) Konvexe Hülle};
  \end{scope}
  % Tendenz
  \draw[ba pfeil] (-0.2,-1.65) -- (13.5,-1.65);
  \node[anchor=north, align=center] at (6.65,-1.75) {von a) nach d) meist engere Anpassung, aber aufwendigerer Überlappungstest};
\end{tikzpicture}
